\section{Research Gap and Problem Formulation}

Phase 1 established the feasibility of a lightweight DOA-sector inference pipeline in which physics-derived beamspace features were evaluated by both a continuous nearest-centroid classifier and a rate-coded LIF-inspired temporal classifier. The completed results showed that the continuous baseline reached an overall accuracy of $86.90\%$, while the best Phase 1 LIF configuration reached $84.17\%$ and exposed measurable stable-decision latency and spike-activity behavior. These results provide a credible reference point, but they do not yet establish whether the same architecture is robust enough to remain useful when the sensing conditions become more difficult.

\subsection{Engineering Research Gap}

The Phase 1 data were generated and evaluated under a controlled single-source ULA model. The training and test samples used the same nominal array geometry, and the reported aggregate accuracy combined samples located throughout each $30^{\circ}$ sector. This design was appropriate for establishing feasibility, but it can conceal two practically important failure mechanisms. First, samples close to a sector boundary are intrinsically more ambiguous than samples near a sector center. Second, a classifier trained with a nominal steering model may degrade when the received snapshots contain element-level gain or phase errors.

A second gap concerns the interpretation of the spike-based path. Phase 1 demonstrated that rate encoding and leaky temporal accumulation can preserve most of the continuous-baseline accuracy, but the result was evaluated primarily under nominal conditions. For an embedded sensing system, accuracy alone is not sufficient. A temporal inference path must also be characterized by the time required to reach a stable decision, the number of generated spike events, and the degree to which those quantities change as signal quality and model fidelity deteriorate.

The Phase 2 research gap is therefore not the absence of another high-resolution DOA estimator. It is the absence of a controlled robustness study that preserves the completed Phase 1 pipeline while comparing classical sector estimation, continuous feature classification, and spike-based temporal inference on identical stressed test conditions. This comparison is needed before the Python reference model can support a defensible transition toward fixed-point analysis, PYNQ deployment, or Vitis HLS implementation.

\subsection{Problem Formulation}

Phase 2 retains the Phase 1 eight-element half-wavelength ULA, sample-covariance calculation, beamspace feature vector, sector definitions, nearest-centroid baseline, rate encoder, and LIF-inspired class accumulator. The nominal four-sector label set remains

\[
\mathcal{C}=\{[-60^{\circ},-30^{\circ}),[-30^{\circ},0^{\circ}),[0^{\circ},30^{\circ}),[30^{\circ},60^{\circ}]\}.
\]

Each Phase 2 test condition is described by a controlled stress vector

\[
\boldsymbol{\eta}
=
\left[
\gamma,
\Delta_{b},
\sigma_{g},
\sigma_{\phi}
\right],
\]

where $\gamma$ is the signal-to-noise ratio, $\Delta_{b}$ is the angular distance from the nearest sector boundary, $\sigma_{g}$ characterizes element-gain variation, and $\sigma_{\phi}$ characterizes element-phase variation. The nominal Phase 1 case corresponds to no array mismatch and unrestricted within-sector angular sampling. Phase 2 will vary these factors systematically rather than replacing the original signal and inference architecture.

For every common test sample, three inference paths will be evaluated:

\begin{enumerate}
    \item a classical array-processing reference whose angular estimate is mapped to one of the four sectors;
    \item the continuous nearest-centroid classifier operating on normalized beamspace features; and
    \item the rate-coded LIF-inspired classifier operating on temporal spike evidence.
\end{enumerate}

The principal outputs will be sector-classification accuracy, confusion behavior, stable-decision latency, and spike activity. To compare the spike-based and continuous paths directly, Phase 2 will also report the accuracy-retention ratio

\[
\rho_{A}
=
\frac{A_{\mathrm{LIF}}}{A_{\mathrm{continuous}}},
\]

where $A_{\mathrm{LIF}}$ and $A_{\mathrm{continuous}}$ are measured on the same test condition. This ratio distinguishes absolute performance loss caused by a difficult sensing condition from the additional loss introduced by temporal spike encoding.

The central research problem is therefore:

\begin{quote}
How robust is the completed Phase 1 spike-based DOA-sector inference pipeline to low SNR, sector-boundary ambiguity, and modest array-model mismatch, and how does its accuracy--latency behavior compare with classical and continuous-valued inference references evaluated on the same data?
\end{quote}

The scope remains sector-level directional inference for autonomous edge sensing. Phase 2 does not claim a new high-resolution DOA algorithm, a complete radar signal chain, or a hardware implementation. Its purpose is to establish the controlled comparative evidence required to justify those later engineering phases.
